% ---------------------------------------------------------------------
% Section 1
% ---------------------------------------------------------------------

\begin{cnread}
本节给出全篇的问题、舞台与记号，并把后面要用到的 \AEG{} 结论一次列清
（引用一律用 Paper~0 / Paper~IV 的\textbf{编译编号}，不引用对话里的口头编号）。
五个「物理字典假设」在这里集中声明，后面每一节用到哪一条都会指明。若只想
看结论，可直接跳第十节的对照表。
\end{cnread}

\subsection{The question}

\AEG{} studies an arithmetic expression through the ordered histories that
evaluate it.  Paper~0 builds three carriers for one history word: the basic
hyperbolic arithmetic expression space $\E_0=(\HH^2,g_{\mu,\lambda},a)$, in
which addition translates and multiplication scales
(\cite{AEG0}, Proposition~3.6 and the table in Subsection~3.4); the
accumulative commutative space $\ACS=\R^2_{(A,M)}$, which records the total
additive charge $A$ and multiplicative charge $M$ and forgets their order
(Definition~8.1); and the three-dimensional contact model
$\Cont=\R^3_{(u,v,a)}$ with
\[
  \alpha = \mathrm da-\mu\,\mathrm du-\lambda a\,\mathrm dv,
  \qquad
  \mathrm d\alpha=-\lambda\,\mathrm da\wedge \mathrm dv,
\]
whose horizontal fields $D_u=\pt_u+\mu\pt_a$ and $D_v=\pt_v+\lambda a\pt_a$
satisfy $[D_u,D_v]=\mu\lambda\,\pt_a$ (Proposition~9.1, Theorem~9.2).  The order
defect of two histories with equal total charges is an exact weighted area in
$\ACS$ (Theorem~8.9); in $\Cont$ the same defect has the constant density
$\mu\lambda$ (Theorem~9.2 and Proposition~9.6); and after passing to the
Reeb quotient $Q(u,v,a)=(-ae^{-\lambda v},e^{-\lambda v})$ of
Proposition~9.8, the defect acquires a second, conjugate reading as a drift
along the Reeb fiber (Theorem~9.9).

This note asks what that machinery says when the two ends of a history are
given a physical reading.  Fix an ideal boundary point $q\in\pt_\infty\HH^2$ and
an ordinary point $p\in\HH^2$.  In \AEG{} terms nothing is new: the histories joining a starting point deep in
the ultraviolet to $p$ inside $B_\infty$ form the fiber of the operator
evaluation map over one affine operator (Proposition~\ref{prop:operator-fiber}).  In physical terms the same object has a standard
name.  The ideal point is an asymptotic state, the ordinary point is a
finite-scale dressed state, the multiplicative steps are scale evolution, and
the additive steps are vertex insertions.  The question is then simply whether
the \AEG{} invariants attached to that family --- the fan-out in $\ACS$, the
Reeb drift, the cross-ratio holonomy, the radial action --- reproduce the
electromagnetic phenomena that the corresponding Feynman histories describe.

\begin{figure}[t]
\centering
\begin{tikzpicture}[x=1cm,y=1cm,>=Latex,every node/.style={font=\small}]
  \node[draw,rounded corners,fill=blue!4,inner sep=4pt,align=center,
        text width=3.6cm] (C) at (0,3.1)
    {$\Cont=\R^3_{(u,v,a)}$\\[1pt]
     \scriptsize$\alpha=\mathrm da-\mu\,\mathrm du-\lambda a\,\mathrm dv$};
  \node[draw,rounded corners,fill=blue!4,inner sep=4pt,align=center] (A) at (-3.5,0.6)
    {$\ACS=\R^2_{(A,M)}$\\[1pt]\scriptsize charges only};
  \node[draw,rounded corners,fill=blue!4,inner sep=4pt,align=center] (H) at (3.5,0.6)
    {$\E_0=\HH^2_{(x,y)}$\\[1pt]\scriptsize assignment $a=-x/y$};
  \draw[->,thick,black!65] (C.south west) -- node[left=1pt,font=\scriptsize]
        {$\pi$} (A.north east);
  \draw[->,thick,black!65] (C.south east) -- node[right=1pt,font=\scriptsize]
        {$Q$} (H.north west);
  \draw[{Latex[length=1.5mm]}-{Latex[length=1.5mm]},dashed,black!45]
        (A.east) -- node[below=2pt,font=\scriptsize]
        {Legendre dual (Thm.~\ref{thm:dual-tearing})} (H.west);
  \node[font=\scriptsize,align=center,text width=13.2cm] at (0,-1.15)
    {$\pi(u,v,a)=(A,M)$ forgets the assignment $a$ and keeps the charges;\quad
     $Q(u,v,a)=(-ae^{-\lambda v},e^{-\lambda v})$ forgets the Reeb coordinate $u$
     and keeps the state.\\[2pt]
     The order defect is a tearing along $\pt_a$ and a drift along $\pt_u$.};
\end{tikzpicture}
\caption[The two projections of the contact model]%
{The two projections of the contact model, and their Legendre conjugacy.  The
charge projection $\pi$ forgets the assignment, so the order defect appears on
its fiber $\pt_a$ as a value tearing; the Reeb quotient $Q$ forgets the additive
charge, so the same defect appears on its fiber $\pt_u$ as a Reeb drift.
Section~\ref{sec:legendre} is the statement that these are two readings of one
contact curvature.}
\label{fig:two-projections}
\end{figure}

\subsection{The \AEG{} statements used, by compiled number}

All citations below are to the compiled numbering of the \AEG{} manuscripts
\cite{AEG0,AEG4}.  The working dialogue from which this note is drawn cited
several of them loosely; where the loose citation named a different statement
than the one actually used, the corrected reference is given here and the fact
is recorded in Appendix~\ref{app:register}.

{\small
\renewcommand{\arraystretch}{1.22}
\begin{longtable}{@{}p{0.215\linewidth}p{0.735\linewidth}@{}}
\toprule
reference & statement used here \\
\midrule
\endfirsthead
\toprule
reference & statement used here \\
\midrule
\endhead
\midrule
\multicolumn{2}{r}{\emph{continued}}\\
\endfoot
\bottomrule
\endlastfoot
Paper~0, Prop.~3.6 & the grid motions act on the assignment by
$a\circ\mathsf X_s=a+s$, $a\circ\mathsf Y_k=ka$. \\
Paper~0, Subsec.~3.4 & the five-row table of motions, including the reflection
$(x,y)\mapsto(-x,y)$ with $a\mapsto-a$. \\
Paper~0, Thm.~3.9 & continuous arithmetic motion:
$\frac{\mathrm d}{\mathrm ds}(a\circ\gamma)=\mu\cos\theta+\lambda(a\circ\gamma)\sin\theta$. \\
Paper~0, Thm.~4.2, Subsec.~4.3 & ripple pencils and ripple histories; the
common tangency point may move with the prefix. \\
Paper~0, Prop.~5.4, Subsec.~5.5 & fixed points of $z\mapsto a+b/z$ solve
$z^2-az-b=0$; general M\"obius fixed points are eigenlines. \\
Paper~0, Props.~7.1, 7.4 & affine cocycle and the left/right Maurer--Cartan
translation forms. \\
Paper~0, Def.~8.1, Prop.~8.2 & the \ACS{} and its weighted evaluation
$\hat\xi=\int_C e^{-M}\mathrm dA$. \\
Paper~0, Thm.~8.6 & only the multiplicative charge descends; the kernel word
$h(p,q)=\cM_q\cA_p\cM_{-q}\cA_{-e^{-q}p}$ has charges
$\bigl(p(1-e^{-q}),0\bigr)$. \\
Paper~0, Thm.~8.9 & weighted torsion--Stokes:
$\tau(\gamma,\delta)=\iint_{\Sigma}e^{M_*-M}\mathrm dA\wedge\mathrm dM$. \\
Paper~0, Props.~9.1, 9.6, 9.8, Thm.~9.2 & the contact form, its curvature
$\mu\lambda$, the Reeb quotient, and $d\alpha=-\mu\lambda Q^*\mathrm d\vol$. \\
Paper~0, Thm.~9.9 & dual tearing and Reeb drift: $\Delta a$ at fixed charges
versus $\Delta A=-\mu\lambda\,\mathrm{Area}_g(\Omega)$ at fixed state. \\
Paper~0, App.~D.1--D.2, D.4 & the oriented rectangle with pulled-back form
$pq\,\mathrm ds\wedge\mathrm dt$; the closed loop
$\ell=(\cA_p,\cM_q,\cA_{-p},\cM_{-q})$ with $\tau=p(1-e^{-q})$; Darboux
coordinates and $R=-(1/\mu)\pt_u$. \\
Paper~IV, Def.~3.1, Thms.~3.2, 3.3, 3.5 & regular bivaluations
$\Biv_K^{\mathrm{reg}}=\PP^1\!\times\!\PP^1\setminus\Delta$; the
bivaluation--projector bijection $\Pi_{a,b}=\frac1{1+ab}v_a\otimes\varphi_b$;
$G/H\cong\Biv_K^{\mathrm{reg}}$; the frame torsor. \\
Paper~IV, Prop.~3.7, Ex.~3.8 & finite-field counts
$|\PGL_2(\F_q)|=q(q^2-1)$, $|\Biv|=q(q+1)$, $|H|=q-1$; for $q=7$:
$336$, $56$, $6$. \\
Paper~IV, Prop.~4.9 & edge labels given only by ratios of frame lifts telescope
along every closed path. \\
Paper~IV, Def.~4.10, Thm.~4.11, Rem.~4.12, Ex.~4.13 & the incidence-coupled
transport $\mathcal T(\tilde e)\in\Aff(1,K)$, its Borel reduction, the
cross-ratio holonomy, the infinitesimal reductive connection, and the
$\F_7$ witness $\Lambda=4$. \\
Paper~IV, Prop.~10.1, Thm.~10.3 & Horner histories
$\xi=\sum_i d_i k^{n-i}$ and their histogram counts. \\
\end{longtable}}

\subsection{Notation and sign conventions}

The note keeps the \AEG{} conventions unchanged; a reader coming from the
manuscripts needs only the following table.

{\small
\renewcommand{\arraystretch}{1.2}%
\begin{longtable}{@{}ll@{}}
\toprule
symbol & meaning \\
\midrule
\endfirsthead
\toprule
symbol & meaning \\
\midrule
\endhead
\midrule
\multicolumn{2}{r}{\emph{continued}}\\
\endfoot
\bottomrule
\endlastfoot
$\E_0=(\HH^2,g_{\mu,\lambda},a)$ & basic hyperbolic arithmetic expression space;
$a(x,y)=-x/y$ \\
$u,v$ & arithmetic direction coordinates; $A=\mu u$, $M=\lambda v$ natural
charge coordinates \\
$a$ & assignment (the ``current arithmetic value''); on $\HH^2$, $a=-x/y$ \\
$\ACS=\R^2_{(A,M)}$ & accumulative commutative space; oriented by
$\mathrm dA\wedge\mathrm dM$ \\
$\Cont=\R^3_{(u,v,a)}$ & contact model; $\alpha=\mathrm da-\mathrm dA-a\,
\mathrm dM$ in charge coordinates \\
$\pi,\;Q$ & charge projection $\pi(u,v,a)=(A,M)$; Reeb quotient
$Q(u,v,a)=(-ae^{-\lambda v},e^{-\lambda v})$ \\
$q,\;p$ & ideal boundary point (asymptote) and ordinary point (finite scale) \\
$\cA_p,\;\cM_q$ & the additive and multiplicative motions, i.e. translation by
$p$ and dilation by $e^{q}$ \\
$h(p,q)$ & the kernel word of Thm.~8.6 \\
$\Pi_{a,b}$ & rank-one idempotent of the bivaluation $(a,b)$ \\
$q^+,q^-$ & image and kernel directions of a bivaluation (Paper~IV) \\
$B_\infty$ & $\Stab_{\PGL_2}(\infty)$, the affine Borel sector \\
$J$ & the inverting element $z\mapsto -1/z$ \\
$\Lambda(\square)$ & multiplicative cross-ratio holonomy of a four-step
bivaluation loop \\
$\mathcal T(e)$ & projective torsion transport of an edge \\
$\ell_0$ & short-distance reference length (ultraviolet cutoff) \\
$\mu,\;\lambda$ & arithmetic step constants; $\mu\lambda$ is the contact
curvature scalar \\
$\kappa$ & interaction curvature coefficient, $\kappa=q_1q_2$ by
Hypothesis~\ref{hyp:pairing} \\
\end{longtable}}

Two conventions matter and are used without further comment.  A history is
written chronologically from left to right, as in Paper~0: the word
$(\cM_q,\cA_p)$ first scales by $e^{q}$ and then translates by $p$.  Charges
are additive as a function on the history group: the total charge of a word is
the sum of its additive and multiplicative increments
(Convention~8.3 of Paper~0).  Composite maps act on the assignment by
composition in the reverse reading, so $(M_k\circ A_h)(a)=e^{q}(a+p)$ while
$(A_h\circ M_k)(a)=e^{q}a+p$; their difference is the elementary order defect
$p(e^{q}-1)$ of Appendix~D.3 of Paper~0.

\subsection{The five dictionary hypotheses}

Everything below is proved either from the \AEG{} statements of
Subsection~1.2 or from elementary computation, except at five points.  Each of
those is declared once here, is used in the sections named, and is not silently
converted into a theorem.

\begin{hypothesis}[Charge--motion dictionary]
\label{hyp:charge}
A charged particle at rest is an ideal-point-to-point history whose additive
steps carry the signed charge: the electron is generated by
$\cA_{-e}$, and the positron by $\cA_{+e}$, which by the reflection of
Paper~0, Subsection~3.4 satisfies $\mathcal R\circ\cA_{-e}\circ\mathcal R^{-1}
=\cA_{+e}$ with $\mathcal R(x,y)=(-x,y)$ scaling $a\mapsto-a$.  A virtual
photon exchange is one additive step of one history absorbed by another.
\end{hypothesis}

\begin{hypothesis}[Depth--distance dictionary]
\label{hyp:depth}
The multiplicative depth $e^{-M}$ of a history is the Coulomb kernel up to the
ultraviolet reference length,
\[
  e^{-M_*}=\frac{\ell_0}{r}
  \qquad\Longleftrightarrow\qquad
  M_*=\log\frac{r}{\ell_0},
\]
so that a scale interval $[0,M_*]$ represents a spatial separation $r$ and the
ultraviolet end $M=0$ represents the coincident limit.
\end{hypothesis}

\begin{hypothesis}[Curvature--energy dictionary]
\label{hyp:pairing}
The two-body interaction energy is the bilinear pairing of the two charges
against the contact curvature.  Concretely, writing the additive generator
strength $\mu\propto q_1$ and the multiplicative coupling $\lambda\propto q_2$,
the curvature scalar of Paper~0, Theorem~9.2 becomes
$\omega(D_u,D_v)=\mu\lambda\propto q_1q_2$, and the potential is the integral of
that density over the scale interval selected by Hypothesis~\ref{hyp:depth},
normalised by the three-dimensional solid angle $4\pi$.
\end{hypothesis}

\begin{hypothesis}[Velocity--tilt dictionary]
\label{hyp:velocity}
A charge moving with speed $v_{\mathrm{vel}}$ is a history whose depth coordinate
advances at rate $\beta=\mathrm dM/\mathrm d\tau=v_{\mathrm{vel}}/c$, i.e. the
direction angle of \eqref{eq:theta-flow} is the rapidity
$\theta=\operatorname{artanh}(v_{\mathrm{vel}}/c)$ of the motion, with $c$ the
critical rate at which the tilt saturates.
\end{hypothesis}

\begin{hypothesis}[Holonomy--radiation dictionary]
\label{hyp:radiation}
A nonvanishing projective holonomy of a closed history loop, in the sense of the
multiplicative component of Paper~IV, Theorem~4.11, is the carrier of
electromagnetic radiation: uniform motion has trivial holonomy and radiates
nothing, and an acceleration that moves the kernel coordinate produces a
holonomy whose magnitude measures radiated power.
\end{hypothesis}

\begin{statusbox}
\textbf{On the status of the hypotheses.}  None of the five is a theorem of
\AEG{}, and none is claimed to be.  Each is a dictionary: it names a physical
quantity by an \AEG{} object.  Hypotheses~\ref{hyp:charge} and
\ref{hyp:depth} fix what the ends and the scale variable of a history mean;
Hypothesis~\ref{hyp:pairing} fixes what an invariant of two histories means;
Hypothesis~\ref{hyp:velocity} fixes what motion means;
Hypothesis~\ref{hyp:radiation} fixes what a holonomy means.  The mathematical
content of every section stands without them; the physical conclusions are
conditional on them, and are stated in that form.  Section~\ref{sec:summary}
collects which conclusion rests on which hypothesis.

One further identification is made later and is not counted among the five
because it is a substitution rather than a dictionary: in
Section~\ref{sec:bohr} the assignment slot of the contact form is filled by the
scale-corrected radial momentum $r\,p_r$ (Lemma~\ref{lem:action-slot}).  It is
flagged there.
\end{statusbox}

\subsection{Roadmap}

{\small
\renewcommand{\arraystretch}{1.22}%
\begin{longtable}{@{}p{0.30\linewidth}p{0.16\linewidth}p{0.46\linewidth}@{}}
\toprule
source note & section & what it contributes \\
\midrule
\endfirsthead
\toprule
source note & section & what it contributes \\
\midrule
\endhead
\midrule
\multicolumn{3}{r}{\emph{continued}}\\
\endfoot
\bottomrule
\endlastfoot
\texttt{费曼图.md} & \S\ref{sec:feynman} & the $p_0\to p$ family as an operator
fiber; \ACS{} fan-out and bare-charge divergence; Reeb-fiber heights \\
\texttt{勒让德形式.md} & \S\ref{sec:legendre} & the exponential Legendre
identity and the dual tearing theorem \\
\texttt{电动力学讨论-01.md} & \S\ref{sec:coulomb} & Coulomb's sign rule from the
sign of $q_1q_2$; the dipole loop; the three geometric readings \\
\texttt{电动力学讨论-02.md} & \S\ref{sec:ampere}, \S\ref{sec:scattering} &
direction angle and the Amp\`ere coefficient; uniform motion in ripple
geometry; M\o ller versus Bhabha sectors \\
\texttt{电动力学讨论-02.md},
\texttt{电动力学讨论-03.md} & \S\ref{sec:radiation},
\S\ref{sec:torsion} & moving poles, cross-ratio holonomy, the torsion transport
and its Borel reduction \\
\texttt{电动力学讨论-03.md} & \S\ref{sec:bohr} & elliptic fixed points, Reeb
holonomy as action, Bohr--Sommerfeld quantization \\
--- & \S\ref{sec:summary} & dictionary table, status ledger, open ends \\
--- & App.~\ref{app:matrices}--\ref{app:register} & exact computations and the
register of verified identities \\
\end{longtable}}
